\section{Sine-on-random blocks}\label{sec:lines}

\subsection{The problem}

Rotating machines (engines, pumps, fans) add tonal lines to the random vibration, at the rotation
frequency and its orders. In a record these lines are often present only in some parts (idle, cruise)
and their frequency drifts with the speed. A line is not a narrow random resonance: its phase is
coherent over many periods, its amplitude is nearly constant, and a sinusoidal response is far less
damaging than a Gaussian response of the same variance (a sine has no rare large cycles: its damage
ratio to a narrow-band Gaussian response of equal variance is $1/\Gamma(1+b/2)$, which is $1/24$ for
$b=8$, from~\eqref{eq:rayleighmoment} with $A = \sqrt2\sigma$ fixed). Treating the lines as part of a
random process therefore misrepresents the damage. MMVSD detects the blocks of the record where lines
stand on the random background (\emph{sine-on-random blocks}), reproduces the lines as deterministic
sines, and designs the random part for the rest.

\subsection{Detection block by block}

The record is cut into blocks of 5\,s. In every block the PSD is estimated with a resolution of
0.5\,Hz (Hann windows, 75\,\% overlap) and the line detector of Section~\ref{sec:sor2} proposes
candidate peaks. A candidate is accepted only if it passes a \emph{phase coherence} test.

\emph{Coherence.} The frequency of the candidate is followed frame by frame inside the block (the
maximum of each short-time spectrum near the candidate, refined by fitting a parabola to the logarithm
of the three largest values), and a straight line $f(t)$ is fitted to these values. The block is then
demodulated along this frequency, $y(t) = x(t)\,e^{-2\pi\ii\int f\,\dd t}$, and low-pass filtered to
$\pm1$\,Hz, which leaves the \emph{complex envelope} $c(t)$ of whatever lies near the candidate. For a
sine $A\cos(\varphi(t)+\psi)$ with $\varphi = 2\pi\int f\dd t$,
\begin{equation}
  x\,e^{-\ii\varphi} = \tfrac{A}{2}e^{\ii\psi} + \tfrac{A}{2}e^{-\ii\psi}e^{-2\ii\varphi},
\end{equation}
and the low-pass filter removes the second term (at twice the line frequency): $c = \tfrac
A2e^{\ii\psi}$ is constant. For narrow-band noise $c(t)$ is a random phasor. With $N$ samples of $c$,
the coherence is
\begin{equation}
  \Phi = \max_{|\delta|\le r/4}\frac{\bigl|\sum_l c_l\,e^{-2\pi\ii\delta t_l}\bigr|^2}{N\sum_l|c_l|^2}\in[0,1],
\end{equation}
where the small residual frequency offset $\delta$ (up to a quarter of the resolution $r$) absorbs
the error of the frequency fit. By the Cauchy--Schwarz inequality $|\sum_lc_l|^2\le N\sum_l|c_l|^2$, with
equality when $c$ is constant: a sine gives $\Phi\approx1$ (slightly less because of the background in
the filter band). Noise of bandwidth $B$ observed for the block duration $T_b$ has a complex envelope
that stays correlated for about $1/B$; the $N$ samples form about $N_{\mathrm{ind}} = BT_b$
independent groups of $N/N_{\mathrm{ind}}$ nearly equal values. Each group sums to about
$(N/N_{\mathrm{ind}})c_g$, the groups add with random phases, so $\E|\sum_lc_l|^2\approx
N_{\mathrm{ind}}(N/N_{\mathrm{ind}})^2\E|c|^2 = N^2\E|c|^2/N_{\mathrm{ind}}$, while $N\sum|c|^2\approx
N^2\E|c|^2$: noise gives $\Phi\approx1/(BT_b)$. A random resonance with $Q = 60$ at 70\,Hz ($B\approx
1.2$\,Hz) over 5\,s gives $\Phi\approx0.17$; a line gives 0.9 or more.

\emph{Tracks.} A candidate with $\Phi\ge0.75$ is a line; candidates of consecutive blocks within 5\,\%
in frequency are linked into \emph{tracks}. A track is kept if at least two of its blocks are strongly
coherent ($\Phi\ge0.75$); between them it may pass through weakly coherent blocks ($\Phi\ge0.5$), so
that a line with a wandering speed or a low signal-to-noise ratio is not cut (hysteresis). A random
resonance that passes the test by chance in one block is dropped.

\subsection{Extraction}

Each track is followed frame by frame in its active intervals (consecutive detected blocks, gaps of up
to two blocks bridged), its frequency is median-filtered over three frames against spurious peaks,
and the record is demodulated along the track's own phase and low-pass filtered (to half the
resolution). By the computation above, the low-pass result is the complex envelope $c(t)$ of the line,
and the line itself is
\begin{equation}
  \ell(t) = 2\,\mathrm{Re}\bigl(c(t)\,e^{\ii\varphi(t)}\bigr),
\end{equation}
tapered at the edges of the intervals. The sum of the lines is subtracted from the record; the
\emph{residual} is the random part of the record. The track keeps its frequency $f(t)$ and amplitude
$a(t) = 2|c(t)|$ on a time grid of the record.

\subsection{The lines of the test}\label{sec:testlines}

The reference time is mapped onto the test time with the mapping of the time--frequency pattern
(Section~\ref{sec:pattern}); with the random field, \code{compress} is used, so that every tonal block
of the record appears in the test, in its order. The amplitude and frequency of each track are read
along the mapped time, smoothed over 0.1\,s, and the track is played as a sine with a random initial
phase.

\emph{Gain of a track.} The damage that a line with a slowly varying amplitude $a(t)$ produces in an
oscillator is the sum of the damages of its cycles, $\frac{K^b}{C}\int f\,(|H|a)^b\,\dd t$. If the
test plays the same frequency path, the factors $f|H|^b$ are the same in the record and in the test,
and the life damage is kept when the test amplitude $g_ja_j$ satisfies
$\int_{\mathrm{test}}(g_ja_j)^b\dd t = \frac{\Tl}{T_{\mathrm{ref}}}\int_{\mathrm{ref}}a_j^b\dd t$:
\begin{result}[Gain of a track]
\begin{equation}
  g_j = \left(\frac{\Tl}{T_{\mathrm{ref}}}\,\frac{\int_{\mathrm{ref}}a_j(t)^b\,\dd t}{\int_{\mathrm{test}}a_j(t)^b\,\dd t}\right)^{1/b}.
\end{equation}
\end{result}
For a stationary line, $\int_{\mathrm{ref}}a^b = A^bT_{\mathrm{ref}}$ and $\int_{\mathrm{test}}a^b = A^bT$,
and $g = (\Tl/T)^{1/b}$, the sine-on-random rule~\eqref{eq:sorgain}. A track whose blocks fall outside
the mapped part of the record is reported; its damage is then carried by the random part.

\subsection{The random part}

Line damage and random damage do not add: a line rides on the random peaks and changes their ranges.
Rather than assuming additivity, MMVSD designs the random part for the \emph{measured} FDS and response
kurtosis of the residual, scaled to the life, $D_{\mathrm{res}}\Tl/T_{\mathrm{ref}}$ and
$\max(\kappa_{\mathrm{res}},3)$, with the PSD of the residual as shape prior and, in pattern mode, the
band-power pattern of the residual. The closed loop then measures the complete test (random part plus
lines) against the complete target, with the lines kept fixed.

\emph{Kurtosis of a response with lines.} Given the field, the response of an oscillator is a
Gaussian part $g$ (variance $s$) plus the responses of the lines, $\sum_j u_j^{1/2}\sqrt2\cos\theta_j$,
where $u_j$ is the response variance of line $j$ and the phases $\theta_j$ are independent and uniformly
distributed. For one line of amplitude $\alpha = \sqrt{2u}$: $\E[\alpha^4\cos^4\theta] =
\alpha^4\cdot\tfrac38 = \tfrac38(2u)^2 = \tfrac32u^2$, using $\cos^4\theta = \tfrac38 +
\tfrac12\cos2\theta + \tfrac18\cos4\theta$. For the sum $X = \sum_jX_j$ of independent zero-mean lines,
the fourth moment keeps the terms $X_j^4$ and the pairs $X_j^2X_l^2$ ($\binom42 = 6$ orderings,
counted once for each unordered pair, i.e.\ $3$ per ordered pair $j\neq l$):
\begin{equation}
  \E[X^4] = \sum_j\tfrac32u_j^2 + 3\sum_{j\ne l}u_ju_l = 3U^2 - \tfrac32\sum_ju_j^2,\qquad U = \sum_ju_j ,
\end{equation}
because $3\sum_{j\ne l}u_ju_l = 3(U^2 - \sum_ju_j^2)$. Adding the independent Gaussian part, whose odd
moments vanish, $\E[(g+X)^4] = \E g^4 + 6\E g^2\E X^2 + \E X^4 = 3s^2 + 6sU + 3U^2 -
\tfrac32\sum u_j^2$:
\begin{result}[Fourth moment of a Gaussian response with lines]
\begin{equation}
  \E[z^4\mid m] = 3(s+U)^2 - \tfrac32\sum_ju_j^2 ,
  \qquad \kappa = \frac{3\tavg{(s+U)^2} - \tfrac32\tavg{\sum_ju_j^2}}{\tavg{s+U}^2}.
\end{equation}
\end{result}
The line variances $u_j(t) = \tfrac12a_j(t)^2|H_d(f_j(t))|^2|H_{\mathrm{pre}}|^4$ are filtered by the
envelope law like $s$. The kurtosis model of the closed loop uses this formula; the damage of the
lines has no reliable closed form on a broad-band random response, so it is not modelled but measured.

\emph{Damage elasticity.} Without lines, multiplying the PSD by $\lambda$ multiplies the amplitudes by
$\sqrt\lambda$ and the damage by $\lambda^{b/2}$: $\partial\ln D/\partial\ln p = b/2$. With lines part of
the damage is fixed and the elasticity $\eta_i$ is lower. It is measured on the first realisation by
amplifying the random part by $1.1$ (its PSD by $1.1^2$):
\begin{equation}
  \eta_i = \frac{\ln\bigl(D_i(1.1\,x_r + x_\ell)/D_i(x_r+x_\ell)\bigr)}{2\ln1.1},
\end{equation}
limited to $[0.1\,b/2,\ b/2]$, and the damage rows of the closed-loop Jacobians are multiplied by
$\eta_i/(b/2)$. Where the lines alone exceed the target damage, the program reports that those
frequencies are over-tested, as the sine-on-random method of version~2.0 does.
